\scanpage{56}
% @01–02, independence
il faut simplement voir que celà ne dépend pas du choix de $\omega^{N-1}$ i.e. que \deleted{$d\omega^{N-1}=0$}
% @03, implication
\[
 d\omega^{N-1}=0\Longrightarrow k_{N-1}(\partial\omega^{N-1})=0
\]
% @04, cases
Distinguons deux cas :

% @05–06, first case
a) $N\geq2$ i.e. $N-1\geq1$, alors $H^{N-1}(A_N^*]=0$ par (1), donc $d\omega^{N-1}=0\Rightarrow\omega^{N-1}=d\alpha^{N-2}$, et on aura
% @07, substitution
\[
 k_{N-1}(\partial\omega^{N-1})=k_{N-1}(\partial d\alpha^{N-2}]\unsure{=|\Leftarrow}
\]
% @08, commutation
Or $\deleted{\partial\alpha^{N-}}\ \partial d\alpha^{N-2}=d\partial\alpha^{N-2}$ donc \unsure{on trouve}
% @09, Stokes
\[
 k_{N-1}(\partial\omega^{N-1})=k_{N-1}(d\partial\alpha^{\unsure{N-1|N-2}})
 \overset{\text{Stokes}}{=}k_{\unsure{N-1|N-2}}(\underbrace{\partial\partial\alpha^{N-2}}_0)=0
\]
% @10, end of case
cqfd

% @11–12, second case
b) $N=1$ i.e. $N-1=0$. Alors en vertu de (1) \deleted{\unsure{on|on a}} $d\omega^{\deleted{\unsure{n}}0}=0$ implique $\omega^{\deleted{\unsure{n}}0}=\varepsilon_1(\alpha)$ ($\alpha\in A_0$) donc
% @13, augmentation
\[
 k_{N-1}(\partial\omega^0)=k_{N-1}(\underbrace{\partial\varepsilon_1(\alpha)}_{=0})=0
 \quad\deleted{\text{or }\partial\varepsilon_1(\alpha)\unsure{=0}}\quad\text{o.k.}
\]
% @14–15, first remark
\uline{Remarques}. 1) Dans \inserted{l'hyp.} (1), en plus de la condition des degrés, on a utilisé
% @16–17, assumptions used
\[
 \left\{\begin{aligned}
 H^N(A_N^*)&=0\quad\text{si }N\geq1\\
 H^{N-1}(A_N^*)&=\begin{cases}0&\text{si }N\geq2\\\varepsilon_1(A_0)&\text{si }N=1\end{cases}
 \end{aligned}\right.
\]
% @18–22, second remark
2) On peut donner une version ``abstraite'' en termes d'un bicomplexe augmenté avec des conditions de degré qu'il faut (et $d\partial=\partial d$ \unsure{!}) \deleted{$\partial\varepsilon_1=0$}) dans une cat. abélienne [\unsure{celà|ceci|celle} peut être $(\mathrm{Ab})^\circ$ !] Conditions
% @18–24, left-hand double complex
\diagram{\begin{tikzcd}[ampersand replacement=\&,column sep=small,row sep=small]
 \&\&\&\&0\&\\
 \&\&\&0\&A_3^3\arrow[u]\&{}\arrow[l]\\
 \&\&0\& A_2^2\arrow[u]\arrow[l]\& A_3^2\arrow[u,"d"']\arrow[l,"\partial"']\&{}\arrow[l]\\
 \&0\& A_1^1\arrow[u]\arrow[l]\& A_2^1\arrow[u,"d"']\arrow[l,"\partial"']\& A_3^1\arrow[u,"d"']\arrow[l,"\partial"']\&{}\arrow[l]\\
 0\& A_0\arrow[l]\& A_1^0\arrow[l,"\partial"']\arrow[u,"d"']\& A_2^0\arrow[l,"\partial"']\arrow[u,"d"']\& A_3^0\arrow[l,"\partial"']\arrow[u,"d"']\&\cdots\arrow[l]\\
 0\arrow[r]\& A_0\arrow[u,"\mathrm{id}","\varepsilon_0"']\& A_0\arrow[l,"0"']\arrow[u,"\varepsilon_1"']\& A_0\arrow[l,"\mathrm{id}"']\arrow[u,"\varepsilon_2"']\& A_0\arrow[l,"0"']\arrow[u,"\varepsilon_3"']\&\cdots\arrow[l,"\mathrm{id}"']
\end{tikzcd}}
% @22–25, revised set of conditions
\deleted{$\unsure{A_N^i=0|A_N^*=0}$, $H^N(A_N^\bullet)=0$ si $N\geq1$}
\[
 \left\{\begin{aligned}
 A_N^i&=0\quad\text{si }i\notin[0,N]\\
 H^N(A_N^*)&=0\quad\text{si }N\geq1\\
 H^{N-1}(A_N^*)&\ \deleted{\unsure{=}}\begin{cases}=0&\text{si }N\geq2\\
 \subset\operatorname{Ker}\delta&\text{si }N=1\end{cases}
 \end{aligned}\right.
\]
% @26–30, third remark
3) A cause de l'unicité, ou grâce à la version ``cat. abélienne'', on voit que les hom ``intégration'' $A_n^n\to A_0^0=A_0$ commutent à tous les endom des $A_*^*$, -- donc respectent une éventuelle structure de $k$-module
% archive sheet number: 55

